% =============================================================================
% 2. ARQUITECTURA DE LA SOLUCION
% =============================================================================
\section{Arquitectura de la soluci\'on}
\label{sec:architecture}

\subsection{Visi\'on general}

El sistema sigue una arquitectura de microservicio stateless donde el cliente env\'ia una imagen y recibe predicciones en formato JSON. Los componentes principales son:

\begin{enumerate}
    \item \textbf{Cliente}: cualquier aplicaci\'on que realice requests HTTP (curl, frontend Angular, Postman).
    \item \textbf{API REST (FastAPI)}: recibe la imagen, la valida, la preprocesa y orquesta la inferencia.
    \item \textbf{Motor de inferencia (ONNX Runtime)}: ejecuta el modelo DenseNet121 exportado.
    \item \textbf{Contenedor Docker}: empaqueta todo en una unidad portable y reproducible.
\end{enumerate}

\subsection{Flujo de datos}

El procesamiento de una solicitud sigue estos pasos:

\begin{enumerate}
    \item El cliente env\'ia un \texttt{POST /predict} con la imagen en formato \texttt{multipart/form-data}.
    \item El \textbf{FilterService} valida que la imagen sea una radiograf\'ia (grayscale) mediante an\'alisis de varianza RGB. Im\'agenes a color son rechazadas.
    \item El \textbf{PreprocessingService} convierte la imagen a escala de grises, la redimensiona a $224 \times 224$ p\'ixeles y la transforma en un tensor float32 de forma $(1, 1, 224, 224)$.
    \item El \textbf{InferenceService} ejecuta el modelo ONNX y obtiene 5 probabilidades independientes.
    \item La API devuelve un JSON con las predicciones por patolog\'ia.
\end{enumerate}

\subsection{Decisiones de dise\~no}

La Tabla~\ref{tab:design-decisions} resume las decisiones t\'ecnicas clave:

\begin{table}[H]
\centering
\caption{Decisiones de dise\~no del servicio.}
\label{tab:design-decisions}
\small
\begin{tabularx}{\textwidth}{lX}
\toprule
\textbf{Decisi\'on} & \textbf{Justificaci\'on} \\
\midrule
FastAPI & Framework async, documentaci\'on autom\'atica (OpenAPI), validaci\'on con Pydantic \\
ONNX Runtime & Reemplaza PyTorch ($\sim$2GB) por un runtime liviano ($\sim$50MB), inferencia en $\sim$20ms \\
uv (gestor de deps) & Resoluci\'on 10--100x m\'as r\'apida que pip, lockfile determinista \\
Multi-stage Docker & Separa build de runtime, reduce tama\~no de imagen \\
Modelo embebido & Arranca con \texttt{docker run} sin vol\'umenes ni configuraci\'on \\
Usuario no-root & Principio de m\'inimo privilegio, requerido en entornos productivos \\
\bottomrule
\end{tabularx}
\end{table}

\subsection{Estructura del proyecto}

\begin{lstlisting}[language={}, caption={Estructura del repositorio chest-xpert-backend.}]
chest-xpert-backend/
+-- .github/workflows/
|   +-- ci.yml              # Lint + Test + Build + Release
|   +-- cd.yml              # Build + Push a GHCR
+-- app/
|   +-- main.py             # App factory + lifespan
|   +-- config.py           # Settings (pydantic-settings)
|   +-- routers/
|   |   +-- health.py       # GET /health
|   |   +-- predict.py      # POST /predict
|   +-- schemas/
|   |   +-- prediction.py   # Request/Response models
|   +-- services/
|       +-- inference.py    # ONNX Runtime session
|       +-- preprocessing.py # Imagen -> tensor
|       +-- filter.py       # Filtro RGB-Diff
+-- models/                  # Modelo ONNX (embebido en imagen)
+-- tests/                   # pytest + Hypothesis
+-- Dockerfile               # Multi-stage production build
+-- pyproject.toml           # Deps + Ruff + Semantic Release
+-- uv.lock                  # Lockfile determinista
\end{lstlisting}
